% Verbatim from /Users/alingalatan/Downloads/cubic_chaos_partial_results (2).zip:reflection_proof.tex (16 November 2026); GPT (a ChatGPT session of the user), relayed by the user.
% Intended filepath: reflection_proof.tex
\documentclass[11pt]{article}
\usepackage[margin=1in]{geometry}
\usepackage{amsmath,amssymb,amsthm,mathtools}
\newtheorem{theorem}{Theorem}
\newtheorem{lemma}{Lemma}
\newtheorem{corollary}{Corollary}
\newcommand{\E}{\mathbb E}
\newcommand{\R}{\mathbb R}
\newcommand{\Hh}{\mathcal H}
\newcommand{\He}{\operatorname{He}}
\newcommand{\norm}[1]{\left\lVert #1\right\rVert}
\title{The Gaussian $L^1/L^2$ constant under reflection oddness}
\author{}
\date{}
\begin{document}
\maketitle

\noindent\textbf{Status and scope.}
Every result below is rigorous.  The theorem proves the Gaussian constant
$\sqrt{2/\pi}$ for pure Gaussian chaoses that are odd under one orthogonal
reflection.  It does \emph{not} prove the bound for arbitrary elements of the
third Wiener chaos: global oddness $h(-x)=-h(x)$ does not imply oddness under a
reflection in a hyperplane.  The proof does not use the extremality assertion
in Example~3.1 of Heilman--Jones--Malavolta.

Let $\He_k$ be the probabilists' Hermite polynomials, normalized by
$\E\He_k(X)^2=k!$ for $X\sim N(0,1)$.  Let $\Hh_d(\R^n)$ denote the span of
the multivariate Hermite polynomials of total degree exactly $d$, with the
Gaussian $L^2$ norm.  Write $c=\sqrt{2/\pi}=\E|X|$.

\begin{theorem}[Reflection-odd chaos]\label{thm:reflection}
Let $d\geq1$ and $h\in\Hh_d(\R^n)$.  Suppose that for a unit vector $e$,
\[
 h\bigl(x-2\langle x,e\rangle e\bigr)=-h(x).
\]
Then
\[
 \E|h|\leq c\norm{h}_2.
\]
More explicitly, after rotating $e$ to the first coordinate and writing
$g=(X,Y)$, there is a unique expansion
\[
 h(X,Y)=\sum_{j=0}^{\lfloor(d-1)/2\rfloor}
     \He_{2j+1}(X)Q_j(Y),\qquad
 Q_j\in\Hh_{d-2j-1}(\R^{n-1}).
\]
With
\[
 b_r=\frac{\binom{2r}{r}}{4^r},\qquad
 A_j=4^j(j!)^2\sum_{r=0}^j b_r^2,
\]
one has the stronger estimate
\begin{equation}\label{eq:strong}
 \E|h|\leq c\left(\sum_j A_j\norm{Q_j}_2^2\right)^{1/2},
 \qquad
 A_j\leq(2j+1)!.
\end{equation}
The coefficient inequality is strict for every $j\geq1$, and $A_0=1$.
\end{theorem}

\begin{proof}
Orthogonal changes of Gaussian coordinates preserve each chaos.  Thus assume
$e$ is the first coordinate vector.  Hermite expansion in $X$ and oddness in
$X$ give the displayed expansion.  Its terms are mutually orthogonal, so
\begin{equation}\label{eq:variance}
 \norm{h}_2^2=\sum_j(2j+1)!\norm{Q_j}_2^2.
\end{equation}

Define a probability measure $\nu$ on the $X$ coordinate by
\[
 d\nu(x)=\frac{|x|}{c}\,d\gamma_1(x),
\]
where $\gamma_1$ is standard Gaussian measure.  Odd Hermite polynomials are
divisible by $x$, so $h(X,Y)/X$ is a polynomial, with its value at zero defined
by continuity.  Cauchy--Schwarz under $\nu\otimes\gamma_{n-1}$ gives
\begin{align}
 \E|h|
 &=c\,\E_{\nu,Y}\left|\frac{h(X,Y)}{X}\right| \notag\\
 &\leq c\left(\E_{\nu,Y}
             \left|\frac{h(X,Y)}{X}\right|^2\right)^{1/2}.\label{eq:bias}
\end{align}
The $Q_j$ belong to distinct Wiener chaoses.  Their orthogonality in $Y$
eliminates all cross terms, even though the $X$ measure has changed.  Hence
\begin{equation}\label{eq:cross}
 \E_{\nu,Y}\left|\frac{h}{X}\right|^2
 =\sum_j\norm{Q_j}_2^2\,
   \E_\nu\left(\frac{\He_{2j+1}(X)}{X}\right)^2.
\end{equation}
It remains to compute the scalar expectations.

Under $\nu$, the variable $T=X^2/2$ has density $e^{-t}$ on $(0,\infty)$.
The Hermite--Laguerre identity is
\begin{equation}\label{eq:hermite-laguerre}
 \frac{\He_{2j+1}(x)}{x}
 =(-1)^j2^j j!L_j^{1/2}(x^2/2).
\end{equation}
For completeness, this follows by matching coefficients in
\[
 \He_m(x)=m!\sum_{k=0}^{\lfloor m/2\rfloor}
   \frac{(-1)^k x^{m-2k}}{2^k k!(m-2k)!},
 \qquad
 L_j^{\alpha}(t)=\sum_{k=0}^j
   \frac{(-1)^k}{k!}\binom{j+\alpha}{j-k}t^k.
\]
The Laguerre generating function
\[
 \sum_{j\geq0}L_j^{\alpha}(t)z^j
 =(1-z)^{-\alpha-1}\exp\left(-\frac{tz}{1-z}\right)
\]
and the binomial series $(1-z)^{-1/2}=\sum_{r\geq0}b_rz^r$ imply, by
coefficient comparison,
\begin{equation}\label{eq:connection}
 L_j^{1/2}(t)=\sum_{r=0}^j b_rL_{j-r}^{0}(t).
\end{equation}
The polynomials $L_k^0$ are orthonormal for $e^{-t}dt$ on $(0,\infty)$.
One can verify this normalization directly from their Rodrigues formula
\[
 L_k^0(t)=\frac{e^t}{k!}\frac{d^k}{dt^k}(e^{-t}t^k):
\]
integration by parts $k$ times gives orthogonality to degrees below $k$;
the leading coefficient $(-1)^k/k!$ then gives squared norm one.
Equations~\eqref{eq:hermite-laguerre}--\eqref{eq:connection} now give
\[
 \E_\nu\left(\frac{\He_{2j+1}(X)}{X}\right)^2
 =4^j(j!)^2\int_0^\infty e^{-t}L_j^{1/2}(t)^2\,dt
 =4^j(j!)^2\sum_{r=0}^j b_r^2=A_j.
\]

Finally, $b_0=1$ and $0<b_r<1$ for $r\geq1$, so $b_r^2\leq b_r$,
strictly when $r\geq1$.  The identity
\begin{equation}\label{eq:sum}
 \sum_{r=0}^j b_r=(2j+1)b_j
\end{equation}
follows by induction from $b_j=(2j-1)b_{j-1}/(2j)$: the difference of the
right-hand sides at $j$ and $j-1$ is exactly $b_j$.
Consequently,
\[
 A_j\leq4^j(j!)^2(2j+1)b_j=(2j+1)!,
\]
with strict inequality for $j\geq1$.  Combining this calculation with
\eqref{eq:bias}, \eqref{eq:cross}, and \eqref{eq:variance} proves all claims.
\end{proof}

\begin{corollary}[Explicit cubic bound]\label{cor:cubic}
If $h\in\Hh_3$ is odd under a hyperplane reflection, then after rotation
\[
 h(X,Y)=a\He_3(X)+XQ(Y),\qquad Q\in\Hh_2(Y),
\]
and
\begin{equation}\label{eq:cubic}
 \E|h|\leq c\sqrt{5a^2+\norm{Q}_2^2}
 \leq c\sqrt{6a^2+\norm{Q}_2^2}=c\norm{h}_2.
\end{equation}
\end{corollary}

\begin{proof}
Theorem~\ref{thm:reflection} gives $A_0=1$ and
$A_1=4(1+1/4)=5$.  There is also a direct proof: with $S=X^2$ under
$\nu$, $S\sim\chi^2_2$, whence $\E_\nu S=2$, $\E_\nu S^2=8$, and
$\E_\nu(S-3)^2=5$.  Since $\E Q=0$ and $X,Y$ remain independent after the
change of measure,
\[
 \E_{\nu,Y}|a(S-3)+Q(Y)|^2=5a^2+\norm{Q}_2^2.
\]
Apply Cauchy--Schwarz as in \eqref{eq:bias}.  Ordinary Hermite orthogonality
gives $\norm{h}_2^2=6a^2+\norm{Q}_2^2$.
\end{proof}

\paragraph{Checks and limitations.}
\begin{itemize}
\item For $d=1$, a nonzero reflection-odd linear form is $h=aX$, and the
bound is an equality: $\E|h|=c\norm{h}_2$.
\item For $d=2$, the reflection-odd case is $h=XL(Y)$ for a linear Gaussian
form $L$ independent of $X$.  Its exact ratio is $c^2=2/\pi<c$ when
$h\ne0$.  This checks the theorem on that subclass; the theorem makes no
claim about \emph{all} quadratic forms based on reflection symmetry.
\item For $h=\He_3(X)$, direct integration gives
\[
 \E|\He_3(X)|=c(1+4e^{-3/2}),\qquad
 \norm{\He_3(X)}_2=\sqrt6.
\]
Indeed, $\E_\nu|S-3|=1+4e^{-3/2}$ for $S\sim\chi^2_2$.
The stronger bound \eqref{eq:cubic} reads
$\E|\He_3(X)|\leq c\sqrt5<c\sqrt6$, consistently.
\item A product of $d$ independent standard Gaussians belongs to $\Hh_d$,
is reflection-odd in each factor, has variance one, and has first absolute
moment $c^d\leq c$.
\item A general cubic has expansion in a chosen direction
\[
 h(X,Y)=a\He_3(X)+(X^2-1)L(Y)+XQ(Y)+R(Y),
\]
with $L\in\Hh_1$, $Q\in\Hh_2$, and $R\in\Hh_3$.
The theorem applies when $L=R=0$ after some rotation.  The other terms
prevent dividing by $X$ while retaining a polynomial with the finite
weighted second moment used in the proof.
\end{itemize}

\end{document}
